% !TEX root = ../DoAn.tex
\documentclass[../DoAn.tex]{subfiles}
\begin{document}

Chương mở đầu này giới thiệu bối cảnh và định hướng tổng thể của đồ án. Trước hết, chương đặt vấn đề về nhu cầu trực quan hóa dữ liệu định vị vệ tinh và chuyển động quang học trên thiết bị di động; tiếp đó xác định mục tiêu, phạm vi và định hướng giải pháp của hệ thống; cuối cùng trình bày phương pháp nghiên cứu và bố cục của toàn bộ đồ án. Các nội dung này lần lượt được trình bày trong các mục Đặt vấn đề (Mục~\ref{section:1.1}), Mục tiêu và phạm vi đề tài (Mục~\ref{section:1.2}), Định hướng giải pháp (Mục~\ref{section:1.3}), Phương pháp nghiên cứu (Mục~\ref{section:1.4}) và Bố cục đồ án (Mục~\ref{section:1.5}).

\section{Đặt vấn đề}
\label{section:1.1}
Trong bài toán định vị trên thiết bị di động, người dùng thường chỉ quan sát kết quả cuối cùng là một điểm tọa độ trên bản đồ. Tuy nhiên, phía sau điểm tọa độ đó là nhiều thông tin kỹ thuật có giá trị, bao gồm số lượng vệ tinh đang được quan sát, cường độ tín hiệu, chòm vệ tinh, trạng thái được dùng trong nghiệm định vị và nguồn dữ liệu quỹ đạo. Android cung cấp các API cho phép truy cập trạng thái vệ tinh và một phần dữ liệu đo GNSS thô, trong đó \texttt{GnssMeasurementsEvent} cho phép ứng dụng nhận tập đo từ GNSS engine trên thiết bị. Nếu được trực quan hóa hợp lý, các dữ liệu này có thể hỗ trợ việc học tập, kiểm thử và đánh giá chất lượng định vị trong các môi trường khác nhau.

Song song với GNSS, camera trên điện thoại có thể cung cấp thêm thông tin về chuyển động tương đối của thiết bị. Optical Flow là nhóm kỹ thuật ước lượng chuyển động biểu kiến giữa các khung hình liên tiếp, có thể dùng để biểu diễn hướng dịch chuyển camera, theo dõi chuyển động trong cảnh hoặc hỗ trợ các bài toán định vị thị giác. OpenCV cung cấp các thuật toán cổ điển như Lucas-Kanade và Farneback, trong khi các mô hình học sâu như RAFT cho phép ước lượng trường chuyển động dày đặc hơn. Mỗi hướng tiếp cận có ưu điểm và chi phí xử lý khác nhau, do đó cần một hệ thống cho phép quan sát, so sánh và đánh giá trong cùng điều kiện sử dụng.

Từ thực tế trên, đồ án lựa chọn xây dựng một hệ thống gồm ứng dụng Android  kết hợp GNSS, camera, Optical Flow và backend xử lý video Optical Flow bằng mô hình mạng lưới học sâu. Hệ thống không nhằm thay thế các ứng dụng bản đồ phổ thông, mà tập trung vào việc trực quan hóa dữ liệu kỹ thuật, thử nghiệm hỗ trợ định vị khi GNSS suy giảm và xem luồng trực quan các thuật toán Optical Flow trong một sản phẩm thử nghiệm thống nhất.

\section{Mục tiêu và phạm vi đề tài}
\label{section:1.2}
\subsubsection{Mục tiêu tổng quát}
Đề tài hướng đến xây dựng một hệ thống phần mềm trên nền tảng Android có khả năng thu nhận, xử lý và trực quan hóa dữ liệu định vị vệ tinh (GNSS) cùng dữ liệu chuyển động quang học (Optical Flow), đồng thời kết hợp một backend nhẹ để xử lý video bằng mô hình mạng lưới học sâu RAFT. Hệ thống ứng dụng các kỹ thuật thị giác máy tính (Computer Vision) trên luồng camera thời gian thực kết hợp cơ chế phân giải vị trí vệ tinh từ nhiều nguồn, nhằm tạo ra một công cụ thử nghiệm trực quan giúp người dùng quan sát, so sánh và đánh giá dữ liệu GNSS và Optical Flow trong cùng một ứng dụng hoàn chỉnh.

\subsubsection{Mục tiêu cụ thể}
Hệ thống đặt ra năm mục tiêu kỹ thuật trọng tâm. Thứ nhất là trực quan hóa dữ liệu GNSS đa nguồn: thu nhận vị trí hiện tại, trạng thái vệ tinh và dữ liệu đo GNSS thô, phân giải vị trí vệ tinh theo chuỗi ưu tiên (PVT thật từ thiết bị, IGS Broadcast Ephemeris, CelesTrak TLE/SGP4 và phương pháp xấp xỉ), rồi hiển thị trên bản đồ hai chiều, mô hình Trái Đất ba chiều và màn hình thực tế tăng cường (AR). Thứ hai là xử lý Optical Flow thời gian thực trên luồng camera bằng hai thuật toán KLT và Farneback, hỗ trợ điều chỉnh độ nhạy, chọn vùng quan tâm (ROI) và bám đối tượng. Thứ ba là cung cấp cơ chế phân tích, so sánh hiệu năng hai thuật toán trên cùng khung hình đầu vào và lưu lại phiên đo các chỉ số như FPS, thời gian xử lý, số điểm hoặc vector hoạt động và độ tin cậy tương đối. Thứ tư là hỗ trợ định vị khi tín hiệu GNSS suy giảm bằng cách kết hợp Optical Flow với cảm biến quán tính (IMU) để ước lượng chuyển động và khớp bản đồ theo tuyến đường hiện có. Cuối cùng là xây dựng một backend nhẹ cho phép tải lên và xử lý video bằng mô hình mạng lưới học sâu RAFT theo cơ chế bất đồng bộ, trả về kết quả chính xác và nhanh hơn so với xử lý trực tiếp trên thiết bị.

\subsubsection{Phạm vi nghiên cứu}
Đề tài tập trung vào việc trực quan hóa dữ liệu kỹ thuật và thử nghiệm các kịch bản kết hợp GNSS, camera và cảm biến quán tính trên thiết bị Android thật, với mẫu giao diện và các luồng xử lý do hệ thống tự xây dựng. Đề tài không hướng đến một hệ thống visual-inertial odometry đầy đủ hay một dịch vụ định vị đạt độ chính xác công nghiệp, mà chỉ dừng ở mức nguyên mẫu kỹ thuật phục vụ quan sát và đánh giá. Ngoài ra, hệ thống chưa đặt mục tiêu triển khai thương mại, chưa tối ưu cho tải nhiều người dùng đồng thời và chưa xây dựng bộ benchmark trên quy mô lớn; phần xử lý video bằng RAFT được đặt ở backend riêng thay vì chạy trực tiếp trên thiết bị.

\section{Định hướng giải pháp}
\label{section:1.3}
\setcounter{subsubsection}{0}
\subsubsection{Kiến trúc client-server và công nghệ nền tảng}
Hệ thống được triển khai theo kiến trúc client-server. Ứng dụng Android đảm nhiệm các tác vụ cần phản hồi trực tiếp với người dùng như thu nhận GNSS, hiển thị bản đồ, dựng mô hình ba chiều và thực tế tăng cường, xử lý Optical Flow thời gian thực và quản lý thư viện kết quả; trong khi các tác vụ xử lý video nặng được chuyển sang backend để tránh phụ thuộc vào tài nguyên của điện thoại. Phía Android sử dụng Kotlin cùng các thư viện CameraX, OpenCV, osmdroid, ARCore, Orekit, WorkManager, Retrofit và Media3; phía máy chủ dùng Python với FastAPI. Cách tách biệt này giúp hệ thống vừa tương tác thời gian thực trên thiết bị, vừa tận dụng được năng lực tính toán của máy chủ.

\subsubsection{Phân giải vị trí vệ tinh đa nguồn theo độ ưu tiên}
Để hiển thị vệ tinh ổn định trên nhiều thiết bị có mức hỗ trợ GNSS khác nhau, hệ thống phân giải vị trí từng vệ tinh theo một chuỗi ưu tiên: ưu tiên PVT thật trích xuất từ thiết bị, sau đó đến IGS Broadcast Ephemeris, kế tiếp là dữ liệu CelesTrak TLE lan truyền bằng SGP4, và cuối cùng là phép xấp xỉ từ azimuth/elevation. Mỗi kết quả được gắn nhãn nguồn dữ liệu để người dùng biết độ tin cậy tương đối, đồng thời bảo đảm giao diện không bị trống khi thiếu nguồn dữ liệu chất lượng cao.

\subsubsection{Xử lý Optical Flow thời gian thực bằng KLT và Farneback}
Trên luồng camera thời gian thực, hệ thống dùng hai thuật toán Optical Flow bổ trợ nhau: KLT cho chuyển động thưa dựa trên điểm đặc trưng và Farneback cho trường chuyển động dày đặc trên toàn ảnh. Người dùng có thể chuyển chế độ hiển thị (vector, heatmap hoặc motion trail), điều chỉnh độ nhạy, khoanh vùng quan tâm (ROI) và bám đối tượng bằng bộ theo dõi CSRT kết hợp template matching. Hệ thống còn cung cấp chế độ phân tích chạy đồng thời hai thuật toán trên cùng khung hình đầu vào để so sánh hiệu năng và lưu lại phiên đo.

\subsubsection{Hỗ trợ định vị khi GNSS suy giảm bằng Optical Flow và IMU}
Khi tín hiệu GNSS suy giảm trong hầm, khu đô thị dày đặc hoặc vùng che khuất, hệ thống bật camera để ước lượng chuyển động bằng Optical Flow, kết hợp tốc độ quay (yaw rate) từ cảm biến quán tính (IMU) và cơ chế ZUPT để khử trôi khi dừng. Vị trí ước lượng được khớp về tuyến đường bằng map matching có điều kiện, đồng thời phân biệt rõ đoạn đường đi bằng GNSS và đoạn được hỗ trợ bằng camera. Đây là một cơ chế hỗ trợ thực dụng, không nhằm thay thế một hệ visual-inertial odometry đầy đủ.

\subsubsection{Xử lý video bằng mô hình mạng lưới học sâu RAFT trên backend}
Các video do người dùng chọn được tải lên backend để xử lý bằng mô hình mạng lưới học sâu RAFT. Máy chủ Python/FastAPI cung cấp các điểm truy cập (endpoint) tạo job, tải lên video, theo dõi trạng thái, hủy job và tải kết quả; pipeline xử lý dùng OpenCV để đọc video, ONNX Runtime để chạy RAFT và ffmpeg để ghi kết quả MP4. Khi người dùng chọn vùng quan tâm, máy chủ có thể dùng Cutie để lan truyền mask đối tượng qua các khung hình. Cơ chế job bất đồng bộ kèm tải lên và tải kết quả theo từng phần (chunk) giúp xử lý video dài mà không phải giữ một kết nối HTTP quá lâu; ngoài máy chủ tự triển khai, ứng dụng còn hỗ trợ gửi video sang một backend ngoài.

\section{Phương pháp nghiên cứu}
\label{section:1.4}
Đề tài được thực hiện bằng cách kết hợp giữa nghiên cứu lý thuyết, khảo sát hiện trạng, thiết kế hệ thống và kiểm thử thực nghiệm. Trước hết, đồ án tìm hiểu các nền tảng kỹ thuật liên quan đến định vị vệ tinh trên Android, cảm biến quán tính, chuyển động quang học, mô hình RAFT và kiến trúc client-server. Đồng thời, một số ứng dụng bản đồ và công cụ Optical Flow hiện có được khảo sát để nhận diện ưu điểm, hạn chế và xác định vấn đề cần giải quyết.Sau đó, hệ thống được triển khai trên thiết bị Android thật, kết hợp với máy chủ xử lý video để kiểm thử trong nhiều điều kiện khác nhau về tín hiệu định vị, ánh sáng và chuyển động camera. Dựa trên kết quả thực nghiệm, các tham số như vùng quan tâm, ngưỡng lọc vector, độ nhạy thuật toán và hệ số hỗ trợ định vị được điều chỉnh nhằm cải thiện độ ổn định của hệ thống. Bên cạnh đó, đồ án sử dụng các mô hình UML như biểu đồ ca sử dụng, tuần tự, hoạt động và lớp để mô tả kiến trúc, luồng xử lý và quan hệ giữa các thành phần.

\section{Bố cục đồ án}
\label{section:1.5}
Đồ án gồm sáu chương. Chương 1 trình bày tổng quan về đề tài, lý do chọn đề tài, mục tiêu, phạm vi và định hướng giải pháp. Chương 2 khảo sát hiện trạng, phân tích yêu cầu, mô tả các ca sử dụng và yêu cầu phi chức năng của hệ thống. Chương 3 trình bày cơ sở lý thuyết và công nghệ sử dụng, bao gồm Android, định vị vệ tinh, hệ tọa độ, bản đồ, mô hình 3D, thực tế tăng cường, Optical Flow, RAFT, ONNX Runtime, FastAPI và cơ chế kết hợp định vị vệ tinh, chuyển động quang học, cảm biến quán tính.Chương 4 mô tả thiết kế và triển khai hệ thống, các mô hình UML, luồng xử lý dữ liệu và kết quả kiểm thử chức năng. Chương 5 phân tích các giải pháp và đóng góp chính của đồ án như phân giải vị trí vệ tinh, hỗ trợ chỉ đường khi tín hiệu định vị suy giảm, đánh giá hiệu năng thuật toán Optical Flow, xử lý video RAFT trên backend và bám đối tượng qua camera. Cuối cùng, Chương 6 tổng kết kết quả đạt được, nêu hạn chế và đề xuất hướng phát triển tiếp theo.

\section*{Kết chương}
Chương này đã làm rõ lý do hình thành đề tài: khoảng trống giữa các ứng dụng bản đồ phổ thông và nhu cầu quan sát, so sánh dữ liệu GNSS cùng chuyển động quang học trong một hệ thống thống nhất. Trên cơ sở đó, đồ án xác định mục tiêu xây dựng ứng dụng Android kết hợp GNSS, camera, Optical Flow và backend xử lý video bằng mô hình mạng lưới học sâu, cùng phạm vi, định hướng giải pháp và phương pháp nghiên cứu tương ứng. Để hiện thực hóa định hướng này, trước hết cần khảo sát hiện trạng và phân tích yêu cầu của hệ thống --- nội dung được trình bày trong Chương~\ref{chapter:Related_works}.

\end{document}
